\input amstex
\documentstyle{amsppt}
\magnification 1200
\pageno 47
\define\ogf{{G(F)}}
\define\ozf{Z_G(F)}
\define\oglf{{GL(2,E)}}
\define\oa{{A_G(F)}}
\define\oslf{{H(E)}}
\define\slf{{H(F)}}
\define\eps{\epsilon}
\define\okk{{\Cal O}_K}
\define\off{{\Cal O}_F}
\define\offp{{\Cal O}_{F'}}
\define\qqq{\Bbb Q}
\define\ppp{\Bbb P}
\define\ccc{\Bbb C}
\define\rrr{\Bbb R}
\define\aaa{\Bbb A}
\define\nnn{\Bbb N}
\define\unt{\underline{T}}
\define\unx{\underline{x}}
\define\oga{{G(\aaa )}}
\define\oba{{B_G(\aaa )}}
\define\oza{{Z_G(\aaa )}}
\define\ogla{{GL(2,\aaa_L )}}
\define\osla{{H(\aaa_L )}}
\define\sla{H(\aaa )}
\define\oee{{\Cal O}_{E}}
\define\oeep{{\Cal O}_{E'}}
\define\VV{\Cal V}
\define\HH{\Cal H}
\hsize 6.5 true in
\vsize 8.7 true in
\vskip .3in
\subheading{\S 7 On the 2-regular germ for $G_2$}
\vskip .1in
Let $G=G_2$ throughtout this section.

The only essential base divisor which can contribute to the 2-regular
germ for $G_2$ via Igusa theory is $E_{\emptyset,\alpha,\delta}
$,
which has $\beta(E_{\emptyset,\alpha,\delta})=3$.
It must be $F$-rational since for every $\sigma\in\Gamma$
$\sigma(E_{\emptyset,\alpha,\delta})$
is also a 2-regular base divisor (see observation (I) in \S 4),
so $\sigma(E_{\emptyset,\alpha,\delta})=E_{\emptyset,\alpha,\delta}$.
Recall that the action of the Galois
group $\Gamma$ on $T_G$ and on $Z$ factors through a finite group
$\Gamma_T$, say.
To satisfy the hypotheses (5.1-4), one may use
the calculations of \S 2 to ``average'' the coordinates of
$Z_{\emptyset,\alpha,\delta}$
over the Galois group $\Gamma_T$
to get $F$-rational coordinates on this divisor if needed.
\par
The branch $E_{\emptyset,\alpha,\delta}$ is defined by
$\lambda =0$ on the regular subset $Z_{\emptyset,\alpha,\delta}$:
$\{x(\delta)\not= 0,w(\delta)\not= 0,z(\beta)\not= 0\}\subset Z$.
The patch $Z_{\emptyset,\alpha,\delta}$ is regular with generators
$z(\alpha)$, $z(\beta)$, $w(\gamma)$, $w(\delta)$, $x(\delta)$,
$x(\epsilon)$, and $x(\zeta)$. The divisor $\lambda =0$ occurs
on this patch as $z(\alpha)=0$ with multiplicity 2. 
The local coordinate relations yield that on $E_{\emptyset,\alpha,\delta}$,
we have $\lambda=x(\alpha)=x(\beta)=z(\alpha)=x(\gamma)=w(\eps)=w(\zeta)=0$.

By \cite{H1, ch 6, \S 2}, only the term $\theta=1$
contributes to the sum over $\theta$ in formula 0.6(a).
The contribution of the 2-regular germ to the $\kappa$-orbital
integral is then
\newline
$\kappa(inv(\gamma_G^0,\gamma_{G,\lambda}))
D_H(\gamma_{H,\lambda})^{-1}$ times
$$
|\lambda|^2
PV\int_{E_{\emptyset,\alpha,\delta}(F)}\pi^*(f)
\kappa_1(t_*)
|\omega_{E_{\emptyset,\alpha,\delta}}|
=|\lambda|^2
PV\int_{{\Cal O}_2}f\int_{E_{\emptyset,\alpha,\delta}(u)(F)}
\kappa_1(t_*)
|\omega_{E_{\emptyset,\alpha,\delta}(u)}|
.
\tag 7.1
$$
\proclaim{Theorem 7.2} Suppose $T_G$ is a non-split Cartan subgroup
of $G_2$ which splits over a
Galois extension for which the Galois group action on $T_G$
is a direct product of two cyclic groups of order $2$.
Assume that $f\in C_c^\infty(G(F))$ satisfies $supp(f)\cap {\Cal O}_2
\not= \emptyset$, where ${\Cal O}_2$ is a fixed $2$-regular
$G(F)$-orbit,
and the support of $f$ does not meet unipotent orbit
of $G(F)$ which is not $2$-regular. Then there is an $f^H$, smooth on
$H(F)$, for which, for all regular $\gamma_G\in T_G(F)$
in a sufficiently small neighborhood of $1$, we have
$$
\Delta_{G/H}^{T_G,\kappa}(\gamma_G)
\Phi^{T_G,\kappa}(\gamma_G,f)=
\Phi^{T_H,st}(\gamma_H,f^H),
$$
where $\kappa$, $T_G$ satisfy
$H_{T_G,\kappa}=(SL(2)\times SL(2))/\{\pm 1\}$. 
(In particular, $\kappa$ is quadratic and $T_G$ is biquadratic.)
\endproclaim                                                         
This will be proven below. 
\vskip .2in 
{\bf Proof of Theorem 7.2}
\vskip .1in
By \cite{Sh} or \cite{Ro} the germ on the
right-hand side is a non-zero constant. The theorem will follow provided
we can show that both sides are equal, as a function of the tangent
direction $X$, up to a constant factor. Indeed, such a constant
factor can be absorbed into
$f^H$ provided we choose $f$ in the beginning so that its support meets the
2-regular orbit and no other unipotent orbit. We fix such a choice 
for the remainder of this section.

We explicitly construct
local coordinates for the unipotent orbit ${\Cal O}_2^{st}$.
Let $u=\theta_{\alpha+\beta }(1)\in G_2$, and compute
the orbit $ad({\frak g}_2)(X_{\alpha+\beta })$.
From this, we conclude that $u^{G_2}$
is 8-dimensional, so ${\Cal O}_2^{st}=u^{G_2}$. 
Analogously to Hales' lemma 6.4,
we find that the coordinates of the subvariety
$\prod_{\sigma\in \Sigma_G-\Sigma_{H_2}}U_\sigma$
of $G_2$ yields local coordinates for the $2$-regular orbit.
Here $\Sigma_{H_2}^+=\{\alpha,\zeta\}$.
There is only one $F$-orbit of $u$ contained in the stable orbit of $u$
(see the appendix), so the $2$-regular orbit of $G_2$ is stable.  
Local coordinates for the fiber $E_{\emptyset,\alpha,\delta}(u)$
may be choosen from the remaining coordinates on
$E_{\emptyset,\alpha,\delta}$:  $z(\beta)$, $x(\zeta)$,
$v(\alpha)$, and $v(\zeta)$. 
\par
We use the form on ${\Cal O}_2$,
$$
\split
\omega_{{\Cal O}_2}
&:={|\bigwedge_{\sigma\in \Sigma_{{\Cal O}_2}}dx(\sigma)\wedge dv(\sigma)|
\over |\lambda|^{3/2}}\\
&={|x(\delta)|^{1/2}\over |w(\delta)|^{3/2}|z(\beta)|^{3/2}}
dz(\alpha)dw(\gamma)dw(\delta)dv(\alpha)dv(\gamma)dv(\delta),
\endsplit
\tag 7.3
$$
where $\Sigma_{{\Cal O}_3}=\Sigma_{G}^+ -\Sigma_{H_2}^+$.
On $E_{\emptyset,\alpha,\delta}(u)$,
$$
\omega_{E_{\emptyset,\alpha,\delta}(u)}:=
{{|x(\delta)|^{5/2}}|w(\delta)|^{7/2}\over{|z(\beta)|^{1/2}}}
dz(\beta)dx(\zeta)dv(\alpha)dv(\zeta).
\tag 7.4
$$
We need to rewrite this last form as an expression
$$
\omega_{E_{\emptyset,\alpha,\delta}(u)}=
|p_{{\Cal O}_2,T_G}(X)|\omega'_{E_{\emptyset,\alpha,\delta}(u)},
$$
where $p_{{\Cal O}_2,T_G}(X)$ is a homogeneous polynomial in the coordinates
of tangent vector $X$ (expressed in a basis of the Lie algebra 
${\frak t}$ of $T_G$) and $\omega'_{E_{\emptyset,\alpha,\delta}(u)}$
is independent of the tangent direction. 
Using (1.22), we find that $p_{{\Cal O}_2,T_G}(X)$ is a product of the form
$$
p_{{\Cal O}_2,T_G}(X)=\prod_{\sigma \in \Sigma_{H_2}^+}\sigma(X)^{f_\sigma},
\tag 7.5
$$
where (by homogeneity of germs)
$\sum_{\sigma >0}f_\sigma =2$. To verify the hypotheses to Hales'
lemma 5.19 one must verify the translation invariance property
of (7.1). The translation invariance of this factor (7.5)
follows from the above expression for
$|p_{{\Cal O}_2,T_G}(X)|$ as a product over roots in $H_2$. 
\par
By continuity, $\kappa(t_*)$ satisfies the 
local translation invariance property as well for sufficiently small $Y$.
Therefore, the hypotheses to Hales' reduction lemma hold and the 
theorem follows in this case from lemma 5.19.

\enddocument



